% ===== 原文 pp.99--102 =====
% [原文 p.99]

注意面积律的势 (4.210) 就好像这两个粒子被一根张力恒定的弦连在一起。这引出了一个著名的想法：禁闭应当与一种动力学上的弦状激发的存在有关，这种激发称为\textbf{禁闭弦}，而两个探针粒子位于它的两端。

\begin{examplebox}{例 4.17：Meissner 效应}
禁闭弦在超导的 Meissner 效应的语境中被具体实现出来。其设定是一个带荷标量场 $\phi$ 的 4d $\U(1)$ 规范理论。当我们给标量场一个非零 vev
\begin{equation}
  \langle\phi\rangle\neq 0\,, \tag{4.211}
\end{equation}
时，与磁荷为 $m$ 的 't Hooft 线算符（它是一个探针 Dirac 磁单极子）相关联的 Dirac 弦就变成一条涡旋弦，使得在其周围的圆周
\begin{equation}
  S^1=\{0\le\theta<2\pi\} \tag{4.212}
\end{equation}
上有
\begin{equation}
  \langle\phi\rangle=v\,\ee^{\ii m\theta}\,. \tag{4.213}
\end{equation}
这是因为沿 Dirac 弦我们有和乐
\begin{equation}
  \oint_{S^1}A=m\,. \tag{4.214}
\end{equation}
这条涡旋弦充当禁闭弦，把 't Hooft 算符禁闭起来。
\end{examplebox}

% [Tutor Note: p.99 (4.213)--(4.214) 的自洽性核对：$\U(1)$ 的和乐 $\exp(\ii\oint A)$ 必须单值，故 (4.214) 中 $m\in\Z$，这正是 't Hooft 线的磁荷量子化。绕 Dirac 弦一圈，$\phi$ 的相位按 (4.213) 变化 $\Delta\arg\phi=2\pi m$，与 (4.214) 给出的 Wilson 线相位一致；两者相乘后场单值。可验证：$m$ 也是标量场凝聚后涡旋弦的绕数。]

\para{库仑行为。}

考虑 IR 有效理论或是纯 Maxwell 理论、或是纯非阿贝尔规范理论的情形。众所周知，这类理论中电探针粒子之间的势随距离按
\begin{equation}
  V(R)\sim\frac{1}{R} \tag{4.215}
\end{equation}
衰减。因此这类理论中的 Wilson 线 $W$ 去禁闭。

磁探针粒子也表现出库仑行为，因此我们断定这类理论中所有双荷线算符都去禁闭。于是我们得到以下陈述：

\begin{statementbox}{陈述 4.9：纯 Maxwell 理论中的 SSB}
$d\ge 4$ 的纯 Maxwell 理论中的电 1-形式对称性与磁（$d-3$）-形式对称性自发破缺。
\end{statementbox}

% [原文 p.100]

\begin{statementbox}{陈述 4.10：纯 Yang--Mills 理论中的 SSB}
$d\ge 5$ 的纯非阿贝尔规范理论中的电 1-形式对称性与磁（$d-3$）-形式对称性自发破缺。
\end{statementbox}

上面两个陈述中时空维数 $d$ 的限制之所以出现，是因为否则规范耦合的 beta 函数为负，从而使规范理论描述不是 IR 中好的有效描述。

关于线算符表现出的其他有趣行为，见 [178]。

\subsubsection{例：具有 $\fsu(2)$ 规范代数的 4d $\cN=1$ SYM 理论}

% [Tutor Note: p.100 小节标题原文写作 “Gauge Algbebra”，是 “Gauge Algebra” 的排版笔误（本节自始至终讨论的是 $\fsu(2)$ 规范代数，且原文别处无此拼写变体）。译文按正确拼写翻译，特此注明以便核对原文。]

让我们用上面这些概念讨论 4d $\cN=1$ $\SU(2)$ 与 $\SO(3)_\pm$ 超对称 Yang--Mills（SYM）理论中的禁闭与去禁闭。

\para{与带软超对称破缺的 $\cN=2$ SYM 的关系。}

这些 $\cN=1$ SYM 理论可以通过取极限
\begin{equation}
  \mu\to\infty \tag{4.216}
\end{equation}
得到：该极限取自纯 $\cN=2$ $\SU(2)$、$\SO(3)_\pm$ SYM 经一个保持 $\cN=1$ 的质量项
\begin{equation}
  \mu\Tr\Phi^2 \tag{4.217}
\end{equation}
形变而后得到的 $\cN=1$ 规范理论。其中 $\Phi$ 是坐落在 $\cN=2$ 矢量多重态之内的 $\cN=1$ 手征多重态。据猜想，$\cN=1$ SYM 理论的相结构与 $\mu$ 取小值的 $\cN=1$ 规范理论的相结构相同。

\para{Seiberg--Witten 理论中的对偶 Meissner 效应。}

Seiberg 与 Witten 在他们那篇里程碑式的论文 [77] 中确定了 $\mu$ 取小值时软形变理论的相结构；他们在其中展示了 Meissner 效应的一个电磁对偶版本，它导致 Wilson 线的禁闭。下面让我们更详细地回顾这个对偶 Meissner 效应。

\para{$\cN=1$ 形变之后的真空。}

$\cN=2$ SYM 理论具有一个真空模空间，称为库仑分支。$\cN=1$ 形变 $\mu$ 会抬升所有这些真空，只有库仑分支上两个特殊的点例外，即磁单极子点与双荷点；它们作为这种软形变之后所得 $\cN=1$ 理论的真空而存活下来。

% [原文 p.101]

\para{IR 有效理论。}

在两个真空中，IR 理论都与上面超导例子中讨论的类似，即一个带规范荷为 1 的标量场 $\phi$ 的 $\U(1)$ 规范理论，该标量场发生凝聚
\begin{equation}
  \langle\phi\rangle\neq 0\,. \tag{4.218}
\end{equation}
此外，我们还有磁荷为偶数的有质量粒子激发。

\para{IR 线算符的禁闭/去禁闭。}

现在让我们研究 IR 理论的各种线缺陷所表现出的行为。首先注意，带电的凝聚体 $\langle\phi\rangle$ 把探针粒子的电荷完全屏蔽掉。因此在计算探针粒子之间的渐近势能时，电荷是无关的。若探针粒子不带任何磁荷，则可以用势 (4.209) 把它们分开而不需任何能量代价，于是所有 Wilson 线都去禁闭
\begin{equation}
  \cW^{IR}_n\longrightarrow\text{去禁闭}\,,\qquad n\in\Z\,. \tag{4.219}
\end{equation}

对于带磁荷的探针粒子，我们需要把偶数磁荷与奇数磁荷的分析分开。带偶数磁荷的动力学粒子的存在，意味着附着在偶数磁荷探针 Dirac 磁单极子上的涡旋弦可以断开，由这些动力学粒子提供新的端点。

因此这类涡旋弦无法产生面积律的势，我们由此得知偶数磁荷的 't Hooft 线去禁闭
\begin{equation}
  \cH^{IR}_{2m}\longrightarrow\text{去禁闭}\,,\qquad m\in\Z\,. \tag{4.220}
\end{equation}

另一方面，奇数磁荷的 't Hooft 线像超导例子中那样禁闭
\begin{equation}
  \cH^{IR}_{2m+1}\longrightarrow\text{禁闭}\,,\qquad m\in\Z\,. \tag{4.221}
\end{equation}

一般而言，对双荷线我们有
\begin{equation}
\begin{aligned}
  \cD^{IR}_{n,2m}&=\cW^{IR}_n\cH^{IR}_{2m}\longrightarrow\text{去禁闭}\,, && n,m\in\Z\,,\\
  \cD^{IR}_{n,2m+1}&=\cW^{IR}_n\cH^{IR}_{2m+1}\longrightarrow\text{禁闭}\,, && n,m\in\Z\,.
\end{aligned}
\tag{4.222}
\end{equation}

\para{UV 与 IR 线算符之间的对应。}

为了把上述禁闭与去禁闭的结果应用到原来的 UV 理论，我们需要知道 UV 与 IR 线算符之间的对应关系，Seiberg 与 Witten 曾详细描述过这一对应。两个真空中的对应都涉及一个电磁对偶，但对偶变换在两个真空中是不同的：

% [Tutor Note: p.101 (4.219)--(4.222) 的判据可逆推自 (4.209)--(4.210)：IR 中 $\cW_n^{IR}$ 对应电探针（无磁荷），其势为常数 (4.209) 故去禁闭；$\cH^{IR}_{2m+1}$ 的涡旋弦不能断开，故面积律禁闭。核对方法：把 $\cD^{IR}_{n,2m}$ 看成电与磁量子数分别为 $n$ 与 $2m$ 的探针，$2m$ 为偶时弦可断（去禁闭），$2m+1$ 为奇时不可断（禁闭），与 (4.220)、(4.221) 一致。]

% [原文 p.102]

\begin{itemize}[leftmargin=1.4em,itemsep=4pt,topsep=1pt]
\item \textbf{磁单极子真空：} 在磁单极子真空中，从 UV 的观点看 IR 标量场 $\phi$ 带磁荷，因此人们说 $\cN=1$ $\SU(2)$ 理论在这个真空中表现出磁单极子凝聚。线算符的映射为
\begin{equation}
\begin{aligned}
  \cW^{UV}_j&\longrightarrow\cH^{IR}_{2j}\,, && j\in\Z/2\,,\\
  \cH^{UV}_n&\longrightarrow\cW^{IR}_n\,, && n\in\Z\,,\\
  \cD^{UV}_{j,n}&\equiv\cW^{UV}_j\cH^{UV}_n\longrightarrow\cD^{IR}_{n,2j}\,, && j\in\Z/2,\ n\in\Z\,,
\end{aligned}
\tag{4.223}
\end{equation}
其中 $\cH^{UV}_n$ 是这样一个 't Hooft 算符：与之相关联的余特征标
\begin{equation}
  \phi:\U(1)\to\SO(3)\,, \tag{4.224}
\end{equation}
在极大环面
\begin{equation}
  \U(1)\subset\SO(3) \tag{4.225}
\end{equation}
上的绕数为 $n$。
\item \textbf{双荷真空：} 在双荷真空中，从 UV 的观点看 IR 标量场 $\phi$ 带双荷，因此人们说 $\cN=1$ $\SU(2)$ 理论在这个真空中表现出双荷凝聚。线算符的映射为
\begin{equation}
\begin{aligned}
  \cW^{UV}_j&\longrightarrow\cH^{IR}_{2j}\,, && j\in\Z/2\,,\\
  \cH^{UV}_n&\longrightarrow\cD^{IR}_{n,-n}\,, && n\in\Z\,,\\
  \cD^{UV}_{j,n}&\equiv\cW^{UV}_j\cH^{UV}_n\longrightarrow\cD^{IR}_{n,2j-n}\,, && j\in\Z/2,\ n\in\Z\,.
\end{aligned}
\tag{4.226}
\end{equation}
\end{itemize}

\para{UV 线算符的禁闭/去禁闭。}

利用上面的映射 (4.223)、(4.226) 以及 IR 线的行为 (4.222)，我们现在可以推出 UV 线所表现出的行为：

\begin{itemize}[leftmargin=1.4em,itemsep=4pt,topsep=1pt]
\item \textbf{磁单极子真空：} 在磁单极子真空中，UV 线表现出以下行为：
\begin{equation}
\begin{aligned}
  \cW^{UV}_j&\longrightarrow\text{禁闭}\,, && j\notin\Z\,,\\
  \cW^{UV}_j&\longrightarrow\text{去禁闭}\,, && j\in\Z\,,\\
  \cH^{UV}_n&\longrightarrow\text{去禁闭}\,, && n\in\Z\,,\\
  \cD^{UV}_{j,n}&\longrightarrow\text{禁闭}\,, && j\notin\Z,\ n\in\Z\,,\\
  \cD^{UV}_{j,n}&\longrightarrow\text{去禁闭}\,, && j\in\Z,\ n\in\Z\,.
\end{aligned}
\tag{4.227}
\end{equation}
\end{itemize}

% [Tutor Note: p.102 (4.227) 可按 (4.223) 与 (4.222) 逐行核对：$\cW^{UV}_j\to\cH^{IR}_{2j}$，再由 (4.220) 知 $2j\in 2\Z$ 即 $j\in\Z$ 时去禁闭、否则禁闭，正是前两行；$\cH^{UV}_n\to\cW^{IR}_n$ 由 (4.219) 恒去禁闭，即第三行；$\cD^{UV}_{j,n}\to\cD^{IR}_{n,2j}$ 取 (4.222) 的下标 $2m=2j$，当 $j\in\Z$ 时去禁闭，对第四、五行。]
